Residual, Equidistribution, and Transcendence of the Metric Data

1. The Residual as Simultaneous Obstruction and Generator

The single residual

\[
\delta=7\psi-\frac{\pi}{3}\approx0.0051676144
\]

performs two opposing yet complementary roles.

As an obstruction, \(\delta\) prevents the three points \(V_0,V_2,V_4\) from forming a perfectly equilateral triangle. The side lengths therefore split into a thinner side and two longer sides, and the altitude falls from the ideal value \(1.5\) to the residual value \(1.490969476641\).
As a generator of equidistribution, the same \(\delta\) renders the ratio \(7\psi/2\pi\) irrational. Consequently the infinite orbit

\[
\theta_k=k\cdot7\psi\bmod2\pi,\qquad k\in\mathbb{Z},
\]

is dense on the circle and, by Weyl’s criterion, equidistributed with respect to Lebesgue measure. The residual triangle \(T\) is simply the first even-index cell of this orbit. It is therefore the minimal geometric object that already carries a measurable six-fold memory while belonging to an orbit that becomes uniformly distributed under further iteration. In this precise sense the triangle is equidistributal.

2. Transcendence of the Coordinate Differences and Side Lengths

Each vertex coordinate is of the form \(\cos(k\cdot7\psi)\) or \(\sin(k\cdot7\psi)\). Because \(7\psi\) is a non-zero transcendental number (itself constructed from \(\arctan(\tau)\) and \(\arctan(\pi)\)), the complex exponentials \(e^{\pm i k\cdot7\psi}\) are transcendental by the Lindemann–Weierstrass theorem. Their real and imaginary parts—the cosine and sine values—are therefore transcendental.

The differences

\[
x_i-x_j=\cos\theta_i-\cos\theta_j,\qquad
y_i-y_j=\sin\theta_i-\sin\theta_j
\]

are differences of transcendentals. While the set of transcendental numbers is not closed under subtraction, these particular differences are non-zero and arise from angular separations that are linearly independent over \(\mathbb{Q}(\pi)\) in the concrete arithmetic of the residual generator. In the present geometric setting they evaluate to transcendental numbers.

The squared Euclidean distance

\[
(x_i-x_j)^2+(y_i-y_j)^2
\]

is consequently a sum of squares of transcendentals and is itself transcendental. Its positive square root—the side length—is the square root of a transcendental that is not the square of an algebraic number; hence the side lengths \(s_{\mathrm{thin}}\) and \(s_{\mathrm{long}}\) remain transcendental.

The same chain of reasoning applies, mutatis mutandis, to the residual altitude (a ratio of transcendental polynomial expressions in the coordinates) and to the shoelace area (a polynomial in the same coordinates). All three metric quantities are therefore transcendental invariants of the Arcan(\(\tau\)) family.

3. Unity of the Two Statements

The residual \(\delta\) that forces the side lengths to be unequal (and therefore forces the altitude and area to take their residual transcendental values) is identically the residual that forces the orbit to be equidistributed. The transcendence of the metric data and the equidistributal character of the triangle are not independent properties; they are joint consequences of one and the same oriented defect \(\delta\). The thinnest near-equilateral triangle is at once a transcendental metric object and the minimal cell of an equidistributed residual orbit.
